\ExamPart{Câu tự luận}{3}{ }

\NQuestion Một xưởng sản xuất dự định làm các hộp quà giống nhau. Mỗi hộp quà gồm đúng $1$ loại bánh, $1$ loại nước uống và $1$ món quà lưu niệm. Biết rằng:
\begin{itemize}[leftmargin=1.5em]
\item Có $6$ loại bánh;
\item Có $4$ loại nước uống;
\item Có $5$ món quà lưu niệm.
\end{itemize}
\begin{SubQuestions}
  \item Hỏi có bao nhiêu cách chọn một hộp quà?
  \item Nếu xưởng muốn tạo mã cho mỗi hộp quà bằng cách ghi theo thứ tự: tên bánh -- nước uống -- quà lưu niệm, thì số mã có thay đổi không? Giải thích.
\end{SubQuestions}
\AnswerLine{$120$}
\SolutionBlock{\begin{SubQuestions}
\item Theo quy tắc nhân, số cách chọn một hộp quà là
\[
6\cdot4\cdot5=120.
\]
\item Mỗi hộp quà được xác định hoàn toàn bởi bộ ba gồm: loại bánh, loại nước uống, món quà lưu niệm. Khi ghi theo thứ tự bánh -- nước uống -- quà lưu niệm thì mỗi cách chọn tương ứng với đúng một mã. Vì vậy số mã vẫn là
\[
120.
\]
\end{SubQuestions}}

\NQuestion Một khu đất hình chữ nhật có chiều dài hơn chiều rộng $8$ m. Người ta làm một lối đi rộng đều $1$ m ở phía trong dọc theo bốn cạnh, nên phần đất còn lại ở giữa là một hình chữ nhật có diện tích $108$ m$^2$. Tính kích thước ban đầu của khu đất.
\AnswerLine{$10\,\text{m}$ và $18\,\text{m}$}
\SolutionBlock{Gọi chiều rộng ban đầu là $x$ m, khi đó chiều dài là $x+8$ m. Sau khi làm lối đi rộng $1$ m xung quanh phía trong, kích thước phần đất còn lại là:
\[
x-2 \quad \text{và} \quad x+8-2=x+6.
\]
Theo đề bài:
\[
(x-2)(x+6)=108.
\]
Khai triển:
\[
x^2+4x-12=108 \Rightarrow x^2+4x-120=0.
\]
Phân tích:
\[
x^2+4x-120=(x+12)(x-10)=0.
\]
Vì $x>0$ nên $x=10$. Do đó chiều dài là
\[
10+8=18.
\]
Vậy kích thước ban đầu của khu đất là $10$ m và $18$ m.}

\NQuestion Trong mặt phẳng tọa độ $Oxy$, cho ba điểm $A(2;1), B(6;3), C(4;-3)$.
\begin{SubQuestions}
  \item Tính tọa độ các vectơ $\overrightarrow{AB}$ và $\overrightarrow{AC}$.
  \item Chứng minh tam giác $ABC$ vuông tại $A$.
  \item Viết phương trình đường thẳng $BC$.
  \item Tính khoảng cách từ điểm $A$ đến đường thẳng $BC$.
\end{SubQuestions}
\AnswerLine{$\overrightarrow{AB}=(4;2),\ \overrightarrow{AC}=(2;-4),\ BC:3x-y-15=0,\ d(A,BC)=\sqrt{10}$}
\SolutionBlock{\begin{SubQuestions}
\item Ta có
\[
\overrightarrow{AB}=(6-2;3-1)=(4;2),\qquad \overrightarrow{AC}=(4-2;-3-1)=(2;-4).
\]
\item Tính tích vô hướng:
\[
\overrightarrow{AB}\cdot\overrightarrow{AC}=4\cdot2+2\cdot(-4)=8-8=0.
\]
Suy ra $\overrightarrow{AB}\perp\overrightarrow{AC}$, nên tam giác $ABC$ vuông tại $A$.
\item Đường thẳng $BC$ đi qua $B(6;3)$ và $C(4;-3)$ có hệ số góc
\[
m=\dfrac{-3-3}{4-6}=3.
\]
Phương trình qua $B$:
\[
y-3=3(x-6)\Leftrightarrow 3x-y-15=0.
\]
\item Khoảng cách từ $A(2;1)$ đến $BC$ là
\[
d(A,BC)=\dfrac{|3\cdot2-1-15|}{\sqrt{3^2+(-1)^2}}=\dfrac{10}{\sqrt{10}}=\sqrt{10}.
\]
\end{SubQuestions}}

\NQuestion Một đội thiện nguyện cần di chuyển từ điểm $M(4;1)$ đến một địa điểm cứu trợ nằm trên đường thẳng
\[
d: x-y+1=0.
\]
\begin{SubQuestions}
  \item Viết phương trình đường thẳng $\Delta$ đi qua $M$ và vuông góc với $d$.
  \item Tìm tọa độ điểm $H$ là giao điểm của $\Delta$ với $d$.
  \item Tính quãng đường ngắn nhất từ $M$ đến đường thẳng $d$.
\end{SubQuestions}
\AnswerLine{$\Delta:x+y-5=0,\ H(2;3),\ MH=2\sqrt2$}
\SolutionBlock{\begin{SubQuestions}
\item Đường thẳng $d:x-y+1=0$ có hệ số góc $1$. Do đó đường thẳng vuông góc với $d$ có hệ số góc $-1$. Qua $M(4;1)$, ta được:
\[
y-1=-(x-4)\Leftrightarrow x+y-5=0.
\]
\item Tọa độ $H$ là nghiệm của hệ
\[
\begin{cases}
x-y+1=0\\
x+y-5=0
\end{cases}
\]
Cộng hai phương trình:
\[
2x-4=0 \Rightarrow x=2.
\]
Suy ra
\[
2-y+1=0 \Rightarrow y=3.
\]
Vậy $H(2;3)$.
\item Quãng đường ngắn nhất từ $M$ đến đường thẳng $d$ là độ dài đoạn vuông góc $MH$:
\[
MH=\sqrt{(4-2)^2+(1-3)^2}=\sqrt{4+4}=2\sqrt2.
\]
\end{SubQuestions}}
